\subsection{Alcance de la implementación}

La integración se realizó sobre TMFQSfullstate, manteniendo inalterada la interfaz lógica del simulador y concentrando los cambios en la capa de
almacenamiento del vector de estado. Esta decisión permite que las operaciones de más alto nivel sigan trabajando con amplitudes complejas, compuertas
y mediciones, mientras que la representación física del estado puede cambiar de densa a comprimida sin afectar la semántica observable del simulador.

En ese contexto, el esquema desarrollado se apoya en tres decisiones concretas. La primera es representar el estado mediante bloques de amplitudes.
La segunda es emplear Blosc2 como biblioteca de compresión sin pérdida para cada bloque. La tercera es introducir una caché LRU de bloques
descomprimidos para amortiguar accesos repetidos durante la aplicación de compuertas. En la implementación actual, esta organización convive con una representación no comprimida de referencia y con una variante alterna basada en ZFP, lo que permite contrastar distintas realizaciones sin modificar el nivel algorítmico del simulador.
